还可以证明, 对于 Hilbert 空间上的紧算子, Theorem~\ref{thm:norm-limit-of-compact-operators} 的逆命题成立.

\begin{Theorem}
\label{thm:compact-hilbert-finite-rank-approximation}
设 $A\in\mathcal{L}(X,H)$ 为紧算子. 那么存在算子序列 $A_n\in\mathcal{L}(X,H)$, 每个算子均具有有限秩, 使得
\begin{equation}\label{eq:compact-hilbert-finite-rank-approximation-limit}
\lim_{n\to\infty}\|A_n-A\|=0.
\end{equation}

\end{Theorem}

\begingroup
\postdisplaypenalty=10000
\renewcommand{\endProof}{\par\nobreak\hfill$\blacksquare$\par\vspace{0.1cm}}
\begin{Proof}
假设 $A$ 不具有有限秩. 由于 $A$ 为紧算子, 其值域是预紧集的可数并, 因而可分. 设 $\{\phi_i\}_{i=1}^\infty$ 为 $\overline{\mathcal{R}(A)}$ 的一组标准正交基, $P_n$ 为从 $\overline{\mathcal{R}(A)}$ 到
\begin{equation}\label{eq:compact-hilbert-finite-dimensional-projection-space}
M_n=\operatorname{span}\{\phi_1,\ldots,\phi_n\},
\end{equation}
的正交投影, 并令 $A_n=P_nA$. 显然 $A_n$ 具有有限秩. 断言 $A_n\to A$. 反设不然, 则 (抽取适当子序列后) 存在序列 $u_n\in X$, 满足 $\|u_n\|=1$ 且 $\|(A-A_n)u_n\|\geq\varepsilon>0$. 抽取子序列后, 可假设 $Au_n$ 收敛到某个极限 $v$. 现在得到
\begin{equation}\label{eq:compact-hilbert-projection-approximation-contradiction}
\begin{aligned}
(A-A_n)u_n&=(I-P_n)Au_n\\
&=(I-P_n)v+(I-P_n)(Au_n-v).
\end{aligned}
\end{equation}
由于该等式右端收敛到零, 左端也收敛到零, 产生矛盾.
\end{Proof}
\endgroup

\begin{Remark}
\label{rem:finite-rank-approximation-banach-space-caveat}
Theorem~\ref{thm:compact-hilbert-finite-rank-approximation} 对一般 Banach 空间不成立. 另一方面, 也不必将像空间限制为 Hilbert 空间. 实际使用的只是存在强收敛到恒等映射的有限维投影. 在应用中重要的大多数 Banach 空间里, 这样的投影确实存在.
\end{Remark}

对于一般 Banach 空间 $X$ 和 $Y$, 可以证明以下结果.

\begin{Theorem}
\label{thm:compact-operator-adjoint-compact}
设 $A\in\mathcal{L}(X,Y)$ 为紧算子. 那么 $A^\times$ 为紧算子.
\end{Theorem}

请读者在 $X$ 和 $Y$ 为 Hilbert 空间这一特殊情形下证明该结果 (Problem~\ref{ex:compact-adjoint-hilbert-proof}).
